\documentclass[a4paper,11pt]{article}
\usepackage[utf8]{inputenc}
\usepackage[T1]{fontenc}
\usepackage[english]{babel}
\usepackage[left=2.5cm,right=2.5cm,top=2cm,bottom=2cm]{geometry}
\usepackage{hyperref}
\usepackage{xcolor}
\usepackage{microtype}
\usepackage{lmodern}

\definecolor{primary}{RGB}{0, 80, 150}
\hypersetup{colorlinks=true, urlcolor=primary}

\begin{document}
\noindent
\textbf{Lo\"ic Aron MBASSI EWOLO} \\
ENSEA -- Double Dipl\^ome Ing\'enieur \\
\href{mailto:loicaron.mbassiewolo@ensea.fr}{loicaron.mbassiewolo@ensea.fr} \quad (+33) 7 59 52 33 98 \\
\href{https://github.com/Nameless0l}{github.com/Nameless0l} \quad \href{https://mbassiloic.tech}{mbassiloic.tech}

\vspace{1em}
\noindent Cergy, le 22 mars 2026

\vspace{1em}
\noindent
\textbf{Objet : Candidature -- Alternance Pilote Projets IA} \\
\textbf{R\'ef\'erence : Covea}

\vspace{1.5em}
\noindent Madame, Monsieur,

Piloter un portefeuille de projets IA en assurance demande vision strat\'egique, rigueur m\'ethodologique et capacit\'e d'orchestration d'\'equipes techniques. Actuellement stagiaire de recherche \`a INRIA Saclay (\'equipe TRIBE, plateau de Saclay), j'am\'eliore ma capacit\'e \`a g\'erer des projets complexes de recherche multidisciplinaires : collaboration avec \'equipes LVMT (\'Ecole des Ponts), coordination avec partenaires du programme national PEPR MOBIDEC, production de rapports scientifiques et jalons it\'eratifs. Cette exp\'erience forge ma capacit\'e \`a structurer des initiatives IA de grande envergure et \'a communiquer les r\'esultats \`a des stakeholders vari\'es.

\vspace{0.8em}

En amont, en tant que Head de la cellule IA de l'\'Ecole Nationale Sup\'erieure Polytechnique de Yaoundé (ENSPY), j'ai men\'e la croissance de cette cellule de plus de 50\%, orchestr\'e des partenariats avec Google DeepMind, et anim\'e des ateliers techniques aupr\`es d'une cinquantaine d'\'etudiants. Parallèlement, j'ai con\c{c}u la plateforme Snappy, syst\`eme de messagerie d'entreprise s\'ecuris\'ee avec architecture Domain-Driven Design : orchestration de 4 modules (API Spring WebFlux, WebSocket Socket.IO, serveur IA Python/Gemini, SDK TypeScript), implantation du protocole Signal pour chiffrement bout en bout, r\'edaction d'une documentation technique d\'etaill\'ee et pr\'esentation des r\'esultats. Ce projet d\'emontre ma capacit\'e \`a conceptualiser des solutions IA complexes et \`a animer les travaux.

\vspace{0.8em}

Auparavant, ma participation \`a plusieurs comp\'etitions IA (1ère place IndabaX 2025 -- IA Sant\'e, 1ère place Mountain Hub 2024, 1ère place CITS 2024) a r\'ev\'el\'e ma capacit\'e \`a identifier les axes d'impact, structurer les travaux et communiquer les r\'esultats aupr\`es de jur\'es composites. Ma formation rigoureuse en math\'ematiques appliqu\'ees et g\'enie logiciel consolide cette vision strat\'egique par une compr\'ehension technique profonde.

\vspace{0.8em}

Passionn\'e par le pilotage de transformations IA responsable en assurance, je suis disponible pour une alternance de 12 mois \`a partir de septembre 2026.

\vspace{1.5em}
\noindent Veuillez agr\'eer, Madame, Monsieur, l'expression de mes salutations distingu\'ees.

\vspace{2em}
\noindent \textbf{Lo\"ic Aron MBASSI EWOLO}

\end{document}